% Generated. Do not edit.
% Projection scriptural-date-location/evidence, version 1: the Scriptural Date and
% Location page, generated whole from research/chronology.toml, the Scripture
% chronology corpus and the source library by tools/tpt document-types refresh.
% document: liturgy/roman-rite/1962/propers/temporal/52-twelfth-after-pentecost
\section*{Appendix: Evidence for the Scriptural Dates}
\label{sec:date-location-evidence}
{\small For each passage, as the Scriptural Date and Location table names it: what is dated, how the corpus weighs each answer, the source's own words, and the sources cited. A subject the corpus does not date is reported as such and given no date.\par}
{\small
\par\smallskip\noindent\textit{Offertory: Ex. 32:11, 13, 14}\par\nobreak
\begin{itemize}[leftmargin=1.2em,itemsep=1pt,topsep=1pt,parsep=0pt]
\item \textbf{Narrated event}. No date: the corpus holds a date for this that it reports from another edition by exception, which is not printed until the corpus records the source's own qualification of it (research pending).
\end{itemize}
\par\smallskip\noindent\textit{Gradual: Ps. 33:2--3 (Vg.; Heb. 34)}\par\nobreak
\begin{itemize}[leftmargin=1.2em,itemsep=1pt,topsep=1pt,parsep=0pt]
\item \textbf{Traditional attribution} (disputed). David's reign: B.C.~1055--1015 (``reigned from 1055 to 1015 B.C.''; The Catholic Encyclopedia, Volume 4 (New York, 1908), article 'David, King,' printed p. 642, opening paragraph). Place: Hebron, then Jerusalem (The Holy Bible, Douay-Rheims (Challoner revision), Project Gutenberg eBook 1581: 2 Kings 5:5).
\item \textbf{Superscription setting} (preferred). David's flight to Achis, king of Geth: A.M.~2944 (Haydock edition reporting Ussher) (``A. M. 2944, A. C. 1060''; Douay Rheims Edition of the Holy Catholic Bible with a Comprehensive Catholic Commentary, William Peter Feeney 2014 Memorial Edition, 1K.21.1-15; Douay Rheims Edition of the Holy Catholic Bible with a Comprehensive Catholic Commentary, William Peter Feeney 2014 Memorial Edition, Ps.70.1). Place: Geth, at the Philistine court (The Holy Bible, Douay-Rheims (Challoner revision), Project Gutenberg eBook 1581: Ps. 33:1, 1 Kings 21:10, 21:15).
\item \textbf{Composition} (preferred). Psalms: Before c.~165 B.C. (``before the Maccabean period, around 165 B.C.; no individual psalm can be dated securely''; Catholic critical chronology; New American Bible, Revised Edition --- USCCB web state (2026), Psalms introduction).
\end{itemize}
\par\smallskip\noindent\textit{Introit: Ps. 69:2--3 (Vg.; Heb. 70), Ps. 69:4 (Vg.; Heb. 70)}\par\nobreak
\begin{itemize}[leftmargin=1.2em,itemsep=1pt,topsep=1pt,parsep=0pt]
\item \textbf{Historical setting}. No date: the chronology corpus has not yet dated this (research pending).
\item \textbf{Composition}. As under Gradual: Ps. 33:2--3 (Vg.; Heb. 34).
\end{itemize}
\par\smallskip\noindent\textit{Alleluia: Ps. 87:2 (Vg.; Heb. 88)}\par\nobreak
\begin{itemize}[leftmargin=1.2em,itemsep=1pt,topsep=1pt,parsep=0pt]
\item \textbf{Historical setting}. As under Introit: Ps. 69:2--3 (Vg.; Heb. 70), Ps. 69:4 (Vg.; Heb. 70).
\item \textbf{Composition} (preferred). The Psalms of the Sons of Korah: Between the days of Isaias and the return from exile (The Catholic Encyclopedia, Volume 12 (1911)).
\end{itemize}
\par\smallskip\noindent\textit{Communion: Ps. 103:13--15 (Vg.; Heb. 104)}\par\nobreak
\begin{itemize}[leftmargin=1.2em,itemsep=1pt,topsep=1pt,parsep=0pt]
\item \textbf{Historical setting}. As under Introit: Ps. 69:2--3 (Vg.; Heb. 70), Ps. 69:4 (Vg.; Heb. 70).
\item \textbf{Composition}. As under Gradual: Ps. 33:2--3 (Vg.; Heb. 34).
\end{itemize}
\par\smallskip\noindent\textit{Gospel: Luke 10:23--37}\par\nobreak
\begin{itemize}[leftmargin=1.2em,itemsep=1pt,topsep=1pt,parsep=0pt]
\item \textbf{Narrated event}. No date: the chronology corpus has not yet dated this (research pending).
\item \textbf{Composition} (disputed). The Gospel of St Luke: c.~A.D.~70 (``About the year 70''; The Catholic Encyclopedia, Volume 14 (New York, 1912)); Before the end of the Roman imprisonment, when the Acts was finished (De auctore, de tempore compositionis et de historica veritate Evangeliorum secundum Marcum et secundum Lucam (1912), responsa-i-ix).
\end{itemize}
\par\smallskip\noindent\textit{Epistle: 2 Cor. 3:4--9}\par\nobreak
\begin{itemize}[leftmargin=1.2em,itemsep=1pt,topsep=1pt,parsep=0pt]
\item \textbf{Composition} (disputed). The Second Epistle to the Corinthians: A.D.~56 (``(1 and 2 Corinthians; Galatians), 56''; The Catholic Encyclopedia, Volume 11 (New York, 1911)); A few months after the First (The Catholic Encyclopedia, Volume 4 (New York, 1908)).
\end{itemize}
\par}
